\chapter{S\'eries temporelles et donn\'ees chronologiques}
\label{ch:temporal-data}

\begin{quote}
\emph{<< Le temps est la variable la plus importante d'une s\'erie temporelle --- et la plus souvent mal trait\'ee. >>}
\end{quote}

% --- hook ---
Le 1\textsuperscript{er} février 2024, un trader de Londres exécute un ordre à <<~01/02/2024~>> que son système américain interprète comme le 2~janvier. Trente millions de livres de perte. Le bug n'est pas dans le modèle, ni dans l'algorithme : c'est un parse de date qui s'est appuyé sur la locale par défaut. Le temps est la variable la plus importante d'une série temporelle, et la plus souvent mal traitée. Resampler des températures par la somme au lieu de la moyenne fait sauter tout votre pipeline. Créer une variable lag-30 sans gérer les NaN en tête de série injecte des bugs muets. Ce chapitre traite ces erreurs comme du contenu pédagogique --- vous les verrez, vous saurez les repérer.
\section{Parser les dates}

La premi\`ere \'etape avec des donn\'ees temporelles est la conversion de cha\^ines en vrais objets datetime. Pandas fournit un parsing robuste, mais les formats r\'eels sont \'etonnamment incoh\'erents.

\begin{minted}{python}
import pandas as pd
import numpy as np

# Jena Climate : mesures m\'et\'eo aux 10 minutes
# T\'el\'echarger : https://www.kaggle.com/datasets/mnassrib/jena-climate
df = pd.read_csv("jena_climate_2009_2016.csv")
print(df.head())
print(f"\nDate column dtype: {df['Date Time'].dtype}")
\end{minted}

\begin{minted}{python}
# Parser le datetime
df["datetime"] = pd.to_datetime(df["Date Time"], format="%d.%m.%Y %H:%M:%S")
df = df.set_index("datetime").drop(columns=["Date Time"])

print(f"Date range: {df.index.min()} to {df.index.max()}")
print(f"Frequency: ~{(df.index[1] - df.index[0])}")
print(f"Total records: {len(df):,}")
\end{minted}

\begin{warningbox}
Ambigu\"it\'e de parsing : << 01/02/2024 >> est-ce le 2 janvier (US) ou le 1\textsuperscript{er} f\'evrier (Europe) ? Sp\'ecifiez toujours le format explicitement avec \texttt{format}. Ne comptez jamais sur l'inf\'erence automatique en production.
\end{warningbox}

\subsection{G\'erer plusieurs formats de date}

\begin{minted}{python}
# Quand une colonne a des formats m\'elang\'es
messy_dates = pd.Series(["2024-01-15", "15/01/2024", "Jan 15, 2024",
                         "15-Jan-2024", "20240115"])

# pd.to_datetime avec mixed g\`ere de nombreux cas
parsed = pd.to_datetime(messy_dates, format="mixed", dayfirst=True)
print(parsed)

# Pour des donn\'ees vraiment d\'esordonn\'ees, parser sur mesure
def parse_flexible(date_str):
    """Essayer plusieurs formats de date."""
    formats = ["%Y-%m-%d", "%d/%m/%Y", "%b %d, %Y",
               "%d-%b-%Y", "%Y%m%d"]
    for fmt in formats:
        try:
            return pd.to_datetime(date_str, format=fmt)
        except (ValueError, TypeError):
            continue
    return pd.NaT
\end{minted}

% ============================================================
\section{R\'e\'echantillonnage}

\begin{minted}{python}
# Downsampling : donn\'ees 10 min vers horaires
hourly = df.resample("h").mean()
print(f"10-min records: {len(df):,}")
print(f"Hourly records: {len(hourly):,}")

# Downsampling vers journalier
daily = df.resample("D").agg({
    "T (degC)": ["mean", "min", "max"],
    "p (mbar)": "mean",
    "rh (%)": "mean",
})
daily.columns = ["temp_mean", "temp_min", "temp_max",
                 "pressure_mean", "humidity_mean"]
print(daily.head())
\end{minted}

\begin{minted}{python}
# Upsampling : journalier vers horaire (avec interpolation)
daily_sample = daily.head(30)
hourly_upsampled = daily_sample.resample("h").interpolate(method="linear")
print(f"Daily records: {len(daily_sample)}")
print(f"Hourly (upsampled): {len(hourly_upsampled)}")
\end{minted}

\begin{preprocesstip}
En downsamplant, r\'efl\'echissez bien \`a la fonction d'agr\'egation. La temp\'erature utilise la moyenne, la pluviom\'etrie utilise la somme, et les min/max utilisent min/max. Une mauvaise agr\'egation produit des r\'esultats absurdes.
\end{preprocesstip}

% ============================================================
\section{Variables d\'ecal\'ees (lag features)}

\begin{minted}{python}
# Variables d\'ecal\'ees : utiliser des valeurs pass\'ees comme pr\'edicteurs
daily["temp_lag1"] = daily["temp_mean"].shift(1)   # hier
daily["temp_lag7"] = daily["temp_mean"].shift(7)   # semaine derni\`ere
daily["temp_lag30"] = daily["temp_mean"].shift(30) # mois dernier

# Variables de variation
daily["temp_change_1d"] = daily["temp_mean"].diff(1)
daily["temp_change_7d"] = daily["temp_mean"].diff(7)

# Variation en pourcentage
daily["temp_pct_change"] = daily["temp_mean"].pct_change()

print(daily[["temp_mean", "temp_lag1", "temp_lag7",
             "temp_change_1d"]].head(10))
\end{minted}

\begin{minted}{python}
# Autocorr\'elation : v\'erifier l'utilit\'e des lags
import matplotlib.pyplot as plt
from pandas.plotting import autocorrelation_plot

fig, axes = plt.subplots(1, 2, figsize=(14, 5))

# Plot d'autocorr\'elation
autocorrelation_plot(daily["temp_mean"].dropna(), ax=axes[0])
axes[0].set_title("Temperature autocorrelation")
axes[0].set_xlim(0, 400)

# Autocorr\'elation partielle (utile pour d\'eterminer l'ordre)
from statsmodels.graphics.tsaplots import plot_pacf
plot_pacf(daily["temp_mean"].dropna(), lags=50, ax=axes[1])
axes[1].set_title("Partial autocorrelation")

plt.tight_layout()
plt.savefig("autocorrelation.png", dpi=150)
plt.show()
\end{minted}

\begin{warningbox}
Les variables d\'ecal\'ees introduisent des NaN en d\'ebut de s\'erie. Une variable lag-30 a 30 valeurs manquantes. Compensez en supprimant les premi\`eres lignes ou en utilisant forward-fill.
\end{warningbox}

% ============================================================
\section{Statistiques mobiles}

\begin{minted}{python}
# Statistiques mobiles : lissent le bruit
daily["temp_rolling7"] = daily["temp_mean"].rolling(window=7).mean()
daily["temp_rolling30"] = daily["temp_mean"].rolling(window=30).mean()
daily["temp_rolling_std7"] = daily["temp_mean"].rolling(window=7).std()

# Moyenne mobile pond\'er\'ee exponentiellement (plus r\'ecent = plus de poids)
daily["temp_ewm7"] = daily["temp_mean"].ewm(span=7).mean()

print(daily[["temp_mean", "temp_rolling7", "temp_rolling30",
             "temp_ewm7"]].tail(10))
\end{minted}

\begin{minted}{python}
# Visualiser les statistiques mobiles
fig, ax = plt.subplots(figsize=(14, 6))

daily["temp_mean"].plot(ax=ax, alpha=0.3, label="Daily mean")
daily["temp_rolling7"].plot(ax=ax, label="7-day rolling mean")
daily["temp_rolling30"].plot(ax=ax, label="30-day rolling mean",
                             linewidth=2)

ax.set_ylabel("Temperature (degC)")
ax.set_title("Jena Climate: Temperature with rolling averages")
ax.legend()
plt.tight_layout()
plt.savefig("rolling_stats.png", dpi=150)
plt.show()
\end{minted}

% ============================================================
\section{D\'ecomposition tendance/saisonnalit\'e}

\begin{minted}{python}
from statsmodels.tsa.seasonal import seasonal_decompose

# D\'ecomposer en tendance + saisonnier + r\'esidu
# Utiliser des donn\'ees mensuelles pour une d\'ecomposition plus propre
monthly = daily["temp_mean"].resample("ME").mean()

decomposition = seasonal_decompose(monthly, model="additive", period=12)

fig = decomposition.plot()
fig.set_size_inches(14, 10)
plt.tight_layout()
plt.savefig("decomposition.png", dpi=150)
plt.show()
\end{minted}

\begin{minted}{python}
# Extraire les composantes comme variables
monthly_df = pd.DataFrame({
    "observed": decomposition.observed,
    "trend": decomposition.trend,
    "seasonal": decomposition.seasonal,
    "residual": decomposition.resid,
})
print(monthly_df.head(15))
\end{minted}

\begin{preprocesstip}
Le r\'esidu de d\'ecomposition est ce qui reste apr\`es retrait de tendance et saisonnalit\'e. Un grand r\'esidu indique une observation inhabituelle --- m\'ecanisme de d\'etection d'outliers pour s\'eries temporelles.
\end{preprocesstip}

% ============================================================
\section{\'Etude de cas Air Quality}

\begin{minted}{python}
# Air Quality UCI : mesures horaires de polluants
df_air = pd.read_csv("AirQualityUCI.csv", sep=";", decimal=",",
                      parse_dates={"datetime": ["Date", "Time"]},
                      dayfirst=True)

# Remplacer les sentinelles
df_air = df_air.replace(-200, np.nan)

# S\'electionner les colonnes cl\'es
cols = ["datetime", "CO(GT)", "NO2(GT)", "T", "RH"]
df_air = df_air[cols].set_index("datetime").dropna(how="all")

# R\'e\'echantillonner en journalier
air_daily = df_air.resample("D").mean()

# Cr\'eer des variables temporelles
air_daily["hour_of_day_peak"] = df_air["CO(GT)"].resample("D").idxmax().dt.hour
air_daily["co_lag1"] = air_daily["CO(GT)"].shift(1)
air_daily["co_rolling7"] = air_daily["CO(GT)"].rolling(7).mean()
air_daily["co_weekend"] = air_daily.index.dayofweek.isin([5, 6]).astype(int)

print(air_daily.head(10))
\end{minted}

% ============================================================
\section{R\'ecapitulatif d'extraction de variables datetime}

\begin{minted}{python}
def extract_datetime_features(df, date_col):
    """Extract a comprehensive set of datetime features."""
    dt = df[date_col] if date_col in df.columns else df.index
    features = pd.DataFrame(index=df.index)

    features["year"] = dt.year
    features["month"] = dt.month
    features["day"] = dt.day
    features["dayofweek"] = dt.dayofweek
    features["hour"] = dt.hour if hasattr(dt, "hour") else 0
    features["quarter"] = dt.quarter
    features["is_weekend"] = dt.dayofweek.isin([5, 6]).astype(int)
    features["day_of_year"] = dt.dayofyear
    features["week_of_year"] = dt.isocalendar().week.astype(int)

    # Encodage cyclique
    features["month_sin"] = np.sin(2 * np.pi * dt.month / 12)
    features["month_cos"] = np.cos(2 * np.pi * dt.month / 12)
    features["dow_sin"] = np.sin(2 * np.pi * dt.dayofweek / 7)
    features["dow_cos"] = np.cos(2 * np.pi * dt.dayofweek / 7)

    return features
\end{minted}

% ============================================================
\section{Exercices}

\begin{exercise}
\begin{enumerate}
  \item Chargez Jena Climate. Parsez le datetime et r\'e\'echantillonnez en journalier. Tracez la temp\'erature moyenne quotidienne pour 2015.
  \item Cr\'eez des variables lag (1, 7, 14, 30 jours) pour la temp\'erature. Calculez la corr\'elation de chaque lag avec la temp\'erature actuelle. Quel lag a la plus forte corr\'elation ?
  \item Appliquez une d\'ecomposition saisonni\`ere sur les donn\'ees CO d'Air Quality. Y a-t-il un motif saisonnier clair ? Que dit la tendance ?
  \item Construisez une fonction compl\`ete d'extraction de variables pour Jena Climate qui cr\'ee : lag, rolling, encodages cycliques et indicateurs de tendance. Appliquez-la et reportez le nombre de variables r\'esultantes.
  \item Avec Air Quality, investiguez si les niveaux de CO diff\`erent significativement entre semaine et week-end. Cr\'eez les visualisations appropri\'ees.
\end{enumerate}
\end{exercise}

% ============================================================
\section{R\'esum\'e du chapitre}

\begin{itemize}
  \item Parsez toujours les dates explicitement avec un format connu ; jamais d'inf\'erence automatique en production.
  \item Le r\'e\'echantillonnage change la r\'esolution temporelle : downsampling agr\`ege, upsampling interpole.
  \item Les variables d\'ecal\'ees capturent les d\'ependances temporelles : la temp\'erature d'hier pr\'edit celle d'aujourd'hui.
  \item Les statistiques mobiles (moyenne, \'ecart-type) lissent le bruit et r\'ev\`elent les tendances.
  \item La d\'ecomposition saisonni\`ere s\'epare tendance, saisonnalit\'e et r\'esidus --- chacun peut servir de variable.
  \item L'encodage cyclique (sinus/cosinus) garantit que d\'ecembre est proche de janvier dans l'espace des variables.
\end{itemize}
